% =============================================================================
% RAPPORT — Application Web de gestion de portefeuille d'investissement
% Projet : GestionPortefeuille (Blazor + EF Core + SQLite)
% Compilation : pdflatex (2 passes) ou latexmk -pdf rapport.tex
% Prérequis : distribution TeX (TeX Live, MiKTeX) avec babel-french
% =============================================================================
\documentclass[11pt,a4paper]{report}

\usepackage[T1]{fontenc}
\usepackage{lmodern}
\usepackage[french]{babel}
\usepackage{microtype}
\usepackage{geometry}
\geometry{margin=2.5cm}
\usepackage{graphicx}
\usepackage{hyperref}
\usepackage{url}
\hypersetup{colorlinks=true, linkcolor=blue, urlcolor=blue, citecolor=blue}
\usepackage{booktabs}
\usepackage{longtable}
\usepackage{array}
\usepackage{enumitem}
\usepackage{amsmath}
\usepackage{listings}
\usepackage{xcolor}
\usepackage{caption}
\usepackage{subcaption}
\usepackage{fancyhdr}
\usepackage{titlesec}

% --- Listings C# (approximatif) ---
\lstdefinelanguage{CSharp}{
  morekeywords={abstract,as,async,await,base,bool,break,byte,case,catch,char,checked,class,const,continue,decimal,default,delegate,do,double,else,enum,event,explicit,extern,false,finally,fixed,float,for,foreach,goto,if,implicit,in,int,interface,internal,is,lock,long,namespace,new,null,object,operator,out,override,params,private,protected,public,readonly,ref,return,sbyte,sealed,short,sizeof,stackalloc,static,string,struct,switch,this,throw,true,try,typeof,uint,ulong,unchecked,unsafe,ushort,using,virtual,void,volatile,while,add,alias,ascending,descending,dynamic,from,get,global,group,into,join,let,orderby,partial,remove,select,set,value,var,where,yield},
  sensitive=true,
  morecomment=[l]{//},
  morecomment=[s]{/*}{*/},
  morestring=[b]",
}
\lstset{
  basicstyle=\ttfamily\small,
  keywordstyle=\color{blue}\bfseries,
  commentstyle=\color{gray}\itshape,
  stringstyle=\color{teal},
  breaklines=true,
  frame=single,
  tabsize=2,
  showstringspaces=false
}

\pagestyle{fancy}
\fancyhf{}
\fancyhead[L]{\small Gestion de portefeuille}
\fancyhead[R]{\small Rapport technique}
\fancyfoot[C]{\thepage}
\renewcommand{\headrulewidth}{0.4pt}

\title{\textbf{Application Web de gestion\\d'un portefeuille d'investissement}\\[0.5em]\large Blazor Server, Entity Framework Core, architecture en couches}
\author{Équipe / Mohamed Jegham & Louay Ben Mansour — Polytech\\Semestre 4 — .NET / C\#}
\date{\today}

\begin{document}

\maketitle
\thispagestyle{empty}

\begin{abstract}
\noindent
Ce document présente la conception et la réalisation d'une application Web complète de \emph{gestion de portefeuille} (actifs financiers, transactions, budget, tableau de bord analytique). L'implémentation repose sur \textbf{.NET~8}, \textbf{Blazor Web App} (mode interactif serveur), \textbf{Entity Framework Core} en approche \emph{Code First} avec base \textbf{SQLite}, et une architecture découpée en couches (\textbf{Core}, \textbf{Infrastructure}, \textbf{Services}, \textbf{Web}). Une extension de type \emph{data science} complète les indicateurs classiques : volatilité des rendements, moyenne mobile (SMA), ratio de Sharpe simplifié, drawdown maximal, backtest de règle de tendance, alertes et graphiques (\emph{Chart.js}). Le rapport détaille les exigences pédagogiques, le modèle de données, les services métier et les pistes d'évolution.
\end{abstract}

\tableofcontents
\clearpage

\chapter{Introduction}

\section{Contexte}

La gestion d'un portefeuille d'investissement (actions, crypto-actifs) nécessite le suivi des positions, des liquidités, des opérations d'achat et de vente, ainsi que des indicateurs de performance. Dans le cadre d'un cours sur le développement d'applications Web avec \emph{.NET}, le présent projet met en œuvre une solution complète répondant à un cahier des charges imposant notamment Blazor, Entity Framework Core, une architecture propre, un tableau de bord analytique et l'utilisation de l'asynchronisme.

\section{Objectifs du projet}

Les objectifs principaux sont les suivants :
\begin{itemize}[leftmargin=*]
  \item proposer une interface Web permettant la \textbf{gestion des actifs} et la \textbf{saisie des transactions} (achat / vente) avec validation ;
  \item suivre un \textbf{budget d'investissement} et l'\textbf{historique des opérations} ;
  \item afficher un \textbf{tableau de bord} avec valeur du portefeuille, gain / perte, répartition par actif et par type, et des graphiques ;
  \item intégrer une \textbf{extension analytique} (créativité) : indicateurs statistiques sur l'historique de prix, aide à la décision, backtest simple ;
  \item respecter une \textbf{architecture en couches} et des bonnes pratiques (séparation UI / services / accès données).
\end{itemize}

\section{Périmètre technique}

Le dépôt logiciel est organisé en plusieurs projets au sein d'une solution \texttt{.slnx} :
\begin{description}
  \item[\texttt{GestionPortefeuille.Core}] Entités, énumérations, modèles de projection (DTO), interfaces de services, options de configuration.
  \item[\texttt{GestionPortefeuille.Infrastructure}] \texttt{DbContext}, migrations EF Core, initialisation des données (\emph{seed}).
  \item[\texttt{GestionPortefeuille.Services}] Implémentations des services (portefeuille, transactions, actifs, analytique, alertes, prix marché).
  \item[\texttt{GestionPortefeuille.Web}] Application Blazor, pages, composants, fichiers statiques, configuration.
\end{description}

La base de données relationnelle est \textbf{SQLite} (fichier \texttt{portfolio.db}), choix adapté à un rendu académique (portabilité, déploiement simple).

\section{Organisation du rapport}

Le chapitre~\ref{ch:exigences} rappelle les exigences du cours et la cartographie fonctionnelle. Le chapitre~\ref{ch:arch} décrit l'architecture logicielle. Le chapitre~\ref{ch:donnees} présente le modèle de données et les migrations. Le chapitre~\ref{ch:real} détaille la réalisation (services, asynchronisme, validation). Le chapitre~\ref{ch:ds} expose l'extension \emph{data science}. Le chapitre~\ref{ch:ui} concerne l'interface utilisateur. Le chapitre~\ref{ch:val} résume la validation. Les annexes fournissent des extraits de code et les commandes utiles.

\chapter{Exigences et fonctionnalités}
\label{ch:exigences}

\section{Exigences pédagogiques (rappel)}

Le projet vise à satisfaire les critères suivants, tels que formulés dans l'énoncé du cours :
\begin{itemize}[leftmargin=*]
  \item \textbf{Blazor Web Application} : utilisation de Blazor en mode interactif serveur (\texttt{InteractiveServer}) pour les pages nécessitant des interactions riches (formulaires, graphiques, thème).
  \item \textbf{Architecture propre} : séparation des responsabilités entre UI, services métier et persistance.
  \item \textbf{Entity Framework Core} : approche \emph{Code First} avec \textbf{migrations} versionnées.
  \item \textbf{CRUD} complet sur les entités principales (actifs, transactions ; budget ajustable).
  \item \textbf{Recherche et filtrage} dynamiques via LINQ (actifs : texte, type ; transactions : pagination et tri).
  \item \textbf{Tableau de bord analytique} : indicateurs, graphiques (répartition, courbe de valeur dans le temps).
  \item \textbf{Validation} : \emph{Data Annotations} sur les entités et formulaires Blazor.
  \item \textbf{Asynchronisme} : méthodes \texttt{async}/\texttt{await} côté services et accès EF (\texttt{ToListAsync}, \texttt{SaveChangesAsync}, etc.).
\end{itemize}

\section{Fonctionnalités métier couvertes}

\begin{table}[h]
  \centering
  \caption{Fonctionnalités par domaine}
  \begin{tabular}{@{}p{4.2cm}p{10cm}@{}}
    \toprule
    \textbf{Domaine} & \textbf{Description} \\
    \midrule
    Actifs & CRUD, symbole unique, type Stock/Crypto, prix courant ; liste paginée, tri, recherche avec temporisation (\emph{debounce}). \\
    Transactions & Saisie achat/vente, calcul du montant total, contrôle de liquidité et de quantité détenue ; historique paginé ; export CSV. \\
    Budget & Consultation et mise à jour du budget initial et de la liquidité disponible. \\
    Portefeuille & Instantané : valeur titres, cash, valeur totale, gain/perte vs budget initial, ROI, répartition par actif et par type. \\
    Historique valeur & Points journaliers (\texttt{PortfolioValuePoint}) mis à jour à la visite du tableau de bord et après transactions. \\
    Marchés (optionnel) & Mode \emph{simulé}, mode \emph{API} (Finnhub + CoinGecko avec cache), import CSV des prix. \\
    Alertes & Seuils configurables : perte vs budget, liquidité basse, volatilité élevée par actif. \\
    Data science & SMA, volatilité annualisée approximative, Sharpe-like, drawdown, signal tendance, backtest SMA. \\
    UX & Thème clair/sombre (Bootstrap), graphiques Chart.js, navigation par menu. \\
    \bottomrule
  \end{tabular}
\end{table}

\section{Extension « créativité »}

L'extension retenue apporte une \textbf{valeur analytique et décisionnelle} complémentaire :
\begin{itemize}[leftmargin=*]
  \item quantification du \textbf{risque} (volatilité, drawdown) ;
  \item \textbf{tendance} par rapport à une moyenne mobile ;
  \item \textbf{comparaison risque/rendement} via un indicateur de type Sharpe simplifié ;
  \item \textbf{simulation historique} d'une règle d'investissement (backtest) pour argumenter à l'oral.
\end{itemize}
Cette extension reste cohérente avec le domaine « portefeuille » et s'appuie sur un historique de prix (généré si absent lors du premier chargement analytique).

\chapter{Architecture logicielle}
\label{ch:arch}

\section{Vue en couches}

L'application suit une architecture en \textbf{quatre projets} alignés sur une \emph{clean architecture} simplifiée :

\begin{enumerate}[leftmargin=*]
  \item \textbf{Core} : aucune dépendance vers Infrastructure ni Web ; contient le modèle conceptuel (entités, enums), les contrats (\texttt{interface}) et les DTO utilisés par l'UI ou les services.
  \item \textbf{Infrastructure} : dépend de Core ; implémente la persistance (EF Core, SQLite), le \texttt{DbContext} et les migrations.
  \item \textbf{Services} : dépend de Core et Infrastructure ; concentre la logique métier, le calcul du portefeuille, les règles de transaction, l'analytique et l'intégration HTTP pour les cours.
  \item \textbf{Web} : point d'entrée ASP.NET Core ; configure l'injection de dépendances, le pipeline HTTP, Blazor et référence les trois autres couches.
\end{enumerate}

\section{Principes appliqués}

\begin{description}
  \item[Dépendances vers l'intérieur] L'UI et le composition root (\texttt{Program.cs}) résolvent les interfaces vers des implémentations concrètes enregistrées dans le conteneur IoC.
  \item[Testabilité] Les services prennent en charge des abstractions (\texttt{DbContext}, \texttt{HttpClient} typé, \texttt{IMemoryCache}, \texttt{IOptions}) injectables.
  \item[Cohésion] Un fichier \texttt{ServiceCollectionExtensions} centralise l'enregistrement des services métier et la configuration des options (\texttt{Alerts}, \texttt{PriceData}).
\end{description}

\section{Flux typique d'une requête utilisateur}

Pour une action sur le tableau de bord :
\begin{enumerate}[leftmargin=*]
  \item le composant Blazor (mode serveur) appelle un service injecté (\texttt{IPortfolioService}, \texttt{IAnalyticsService}, \texttt{IPortfolioAlertService}) ;
  \item le service interroge \texttt{AppDbContext} de manière asynchrone ;
  \item les résultats sont projetés en modèles (\texttt{PortfolioSnapshot}, \texttt{AssetAnalytics}, \texttt{PortfolioAlert}) ;
  \item l'UI met à jour l'état et, le cas échéant, invoque du JavaScript (Chart.js) pour le rendu graphique.
\end{enumerate}

\section{Diagramme de classes}

Un diagramme de classes Mermaid est fourni dans le dépôt à la racine : \texttt{DiagrammeClasses.md}. Il est recommandé d'\textbf{exporter} les figures (PNG/SVG) depuis \url{https://mermaid.live} et de les insérer ici avec :

\begin{verbatim}
\begin{figure}[h]
  \centering
  \includegraphics[width=\textwidth]{figures/domaine-ef.png}
  \caption{Modèle entités-associations (extrait).}
\end{figure}
\end{verbatim}

Créer un dossier \texttt{Rapport/figures/} et y placer les exports pour que la compilation \LaTeX{} réussisse.

\chapter{Modèle de données et persistance}
\label{ch:donnees}

\section{Entités principales}

\begin{description}
  \item[\texttt{Asset}] Représente un instrument : symbole (unique), nom, type (\texttt{Stock} / \texttt{Crypto}), prix courant. Collections de navigation vers les transactions et l'historique de prix.
  \item[\texttt{Transaction}] Mouvement d'achat ou de vente : quantité, prix unitaire, date, montant total dérivé ; clé étrangère vers \texttt{Asset}.
  \item[\texttt{PriceHistory}] Série temporelle de clôtures simulées ou issues d'une évolution de marché ; contrainte d'unicité \((\texttt{AssetId}, \texttt{Date})\).
  \item[\texttt{PortfolioBudget}] Budget initial et liquidité disponible (mise à jour lors des opérations).
  \item[\texttt{PortfolioValuePoint}] Valeur totale du portefeuille (cash + positions au cours du jour) pour une date donnée ; index unique sur la date.
\end{description}

\section{Relations et intégrité}

Les relations \texttt{Asset} $\rightarrow$ \texttt{Transaction} et \texttt{Asset} $\rightarrow$ \texttt{PriceHistory} sont en \textbf{one-to-many} avec suppression en cascade côté dépendant, configurées dans \texttt{OnModelCreating}. L'index unique sur \texttt{Asset.Symbol} garantit l'absence de doublons de tickers.

\section{Migrations}

Les migrations EF Core sont stockées dans \texttt{GestionPortefeuille.Infrastructure/Migrations/}, notamment :
\begin{itemize}[leftmargin=*]
  \item création initiale des tables (\texttt{InitialCreate}) ;
  \item ajout de \texttt{PortfolioValuePoints} (\texttt{AddPortfolioValuePoint}).
\end{itemize}

Au démarrage de l'application, \texttt{Database.Migrate()} applique automatiquement les migrations en attente, ce qui facilite la démonstration sur une machine neuve. Un \emph{seed} insère un budget par défaut et des actifs de démonstration si la base est vide.

\section{Choix SQLite}

SQLite convient à un projet étudiant : fichier unique, pas de serveur à installer, sauvegarde du dépôt possible en excluant \texttt{*.db} via \texttt{.gitignore}. Pour une mise en production, une base serveur (PostgreSQL, SQL Server) serait préférable (concurrence, sauvegardes, taille).

\chapter{Réalisation technique}
\label{ch:real}

\section{Composition racine et pipeline}

Le fichier \texttt{Program.cs} de l'application Web enregistre Blazor (composants interactifs serveur), le \texttt{DbContext} SQLite, et les services métier via une méthode d'extension :

\begin{lstlisting}[language=CSharp, caption={Extrait typique d'enregistrement des services}]
builder.Services.AddDbContext<AppDbContext>(options =>
    options.UseSqlite(builder.Configuration.GetConnectionString("DefaultConnection")));
builder.Services.AddPortfolioServices(builder.Configuration);
\end{lstlisting}

Après construction du pipeline (fichiers statiques, antiforgery, Razor Components), une portée \texttt{scoped} exécute les migrations et le \emph{seed}.

\section{Services métier}

\subsection{Portefeuille (\texttt{PortfolioService})}

Calcule l'instantané \texttt{PortfolioSnapshot} à partir des actifs, des transactions et du budget. Le \textbf{coût moyen pondéré} (CMP) des titres encore détenus est obtenu par une passe chronologique sur les opérations (\texttt{PositionCostHelper}) : à l'achat, le CMP est recalculé ; à la vente partielle, le CMP unitaire des titres restants est conservé (logique courante pour une valorisation de portefeuille pédagogique).

\subsection{Transactions (\texttt{TransactionService})}

Vérifie la liquidité pour les achats et la quantité détenue pour les ventes, met à jour \texttt{PortfolioBudget.AvailableCash}, persiste la transaction, puis enregistre un point d'historique de valeur totale via \texttt{IPortfolioService}.

\subsection{Actifs (\texttt{AssetService})}

CRUD, filtrage LINQ (recherche insensible à la casse sur symbole/nom, filtre par type), pagination serveur (\texttt{PagedResult<T>}).

\subsection{Prix marché (\texttt{PriceDataService})}

Selon \texttt{PriceData:Mode} (\texttt{Simulated}, \texttt{Api}, \texttt{None}), met à jour ou non les \texttt{CurrentPrice}. En mode API, utilisation de Finnhub (actions, clé API) et CoinGecko (cryptos mappées), avec mise en cache mémoire pour limiter les appels.

\subsection{Alertes (\texttt{PortfolioAlertService})}

Agrège instantané et analytique pour produire une liste de \texttt{PortfolioAlert} selon les seuils \texttt{appsettings.json}.

\section{Validation}

Les entités \texttt{Asset}, \texttt{Transaction}, \texttt{PortfolioBudget}, etc. portent des attributs \texttt{[Required]}, \texttt{[Range]}, \texttt{[StringLength]}. Les formulaires Blazor utilisent \texttt{EditForm}, \texttt{DataAnnotationsValidator} et \texttt{ValidationSummary}.

\section{Asynchronisme}

Toutes les opérations d'accès base dans les services sont asynchrones (\texttt{Task}, \texttt{async}/\texttt{await}), ce qui évite de bloquer les threads du pool lors des E/S réseau et base de données.

\chapter{Extension analytique et data science}
\label{ch:ds}

\section{Données d'entrée}

L'analytique s'appuie sur \texttt{PriceHistory} : si aucun historique n'existe pour un actif, une série de prix est \textbf{simulée} (marche aléatoire bornée, graine dérivée du symbole) sur une fenêtre de jours configurable, afin de garantir des graphiques et des métriques en démonstration sans dépendre d'Internet.

\section{Moyenne mobile simple (SMA)}

Pour une fenêtre de $N$ jours (ex.\ $N=20$), la SMA au jour $t$ est la moyenne arithmétique des $N$ derniers prix de clôture :
\begin{equation}
  \mathrm{SMA}_N(t) = \frac{1}{N}\sum_{i=0}^{N-1} P_{t-i}.
\end{equation}
Un \textbf{signal de tendance} grossier compare le prix courant à la SMA (écart en pourcentage) : au-dessus (hausse), en dessous (baisse), zone neutre autour de zéro.

\section{Volatilité}

Les rendements journaliers sont définis par $r_t = \dfrac{P_t - P_{t-1}}{P_{t-1}}$. L'écart-type empirique $\hat{\sigma}$ des $r_t$ sur l'échantillon disponible est calculé (estimateur non biaisé avec dénominateur $n-1$). Une \textbf{volatilité annualisée approximative} est obtenue par
\begin{equation}
  \sigma_{\mathrm{ann}} \approx \hat{\sigma} \times \sqrt{252},
\end{equation}
en supposant $252$ jours de négociation par an (hypothèse standard en finance, ici à visée pédagogique).

\section{Ratio de Sharpe simplifié}

Sans taux sans risque explicite ($r_f = 0$), un indicateur \emph{Sharpe-like} annualisé est :
\begin{equation}
  S \approx \frac{\bar{r}}{\hat{\sigma}} \times \sqrt{252},
\end{equation}
où $\bar{r}$ est la moyenne des rendements journaliers. Ce rapport mesure le rendement moyen par unité de risque (écart-type) ; il doit être interprété avec prudence sur des séries courtes ou simulées.

\section{Drawdown maximal}

Sur la série de prix $(P_t)$, on suit le \emph{peak} courant et l'écart relatif au peak. Le drawdown au temps $t$ est $\dfrac{P_t - \mathrm{peak}_t}{\mathrm{peak}_t}$. Le \textbf{maximum drawdown} est le minimum (le plus négatif) de ces rapports, exprimé en valeur absolue en pourcentage pour l'affichage.

\section{Backtest de règle SMA}

À partir du jour où la SMA est disponible, une stratégie simplifiée :
\begin{itemize}[leftmargin=*]
  \item si le prix dépasse la SMA et qu'aucune position n'est détenue, investir la totalité du cash au prix du jour ;
  \item si le prix passe sous la SMA et qu'une position est ouverte, tout liquider au prix du jour.
\end{itemize}
Le résultat est comparé à un \textbf{buy-and-hold} : investissement intégral au premier jour de la fenêtre utile, conservation jusqu'à la fin. Le nombre d'opérations de la stratégie est comptabilisé. Ce backtest illustre l'impact d'une règle technique sur un historique donné (sans frais, sans slippage --- limites à mentionner à l'oral).

\section{Intérêt pour la note « créativité »}

Ces éléments relient explicitement le projet à des notions de \textbf{statistique descriptive}, de \textbf{mesure du risque} et de \textbf{simulation décisionnelle}, tout en restant dans le périmètre du cours (.NET, LINQ, affichage Blazor).

\chapter{Interface utilisateur (Blazor)}
\label{ch:ui}

\section{Structure de navigation}

L'application propose un menu latéral avec les entrées : tableau de bord, actifs, transactions, budget, outils (marchés, import CSV, backtest). Le layout principal inclut un commutateur de \textbf{thème clair/sombre} (attribut \texttt{data-bs-theme} et stockage \texttt{localStorage}).

\section{Pages principales}

\begin{description}
  \item[Tableau de bord] Cartes KPI (valeur totale, liquidité, gain/perte, ROI), alertes, graphiques Chart.js (répartition en anneau, courbe d'historique de valeur), répartition par actif et par type, bloc analytique par symbole, tableau de performance avec CMP et coût de revient.
  \item[Actifs] Formulaire CRUD, filtres, pagination, tri, recherche avec \emph{debounce}.
  \item[Transactions] Saisie avec validation, résumé budget, historique paginé, tri, export CSV global.
  \item[Budget] Ajustement du budget initial et de la liquidité.
  \item[Outils] Documentation des modes de prix, bouton de rafraîchissement API, zone de texte pour import CSV, formulaire de backtest.
\end{description}

\section{Technologies front}

\begin{itemize}[leftmargin=*]
  \item \textbf{Bootstrap} (fichiers locaux sous \texttt{wwwroot/lib/bootstrap}) pour la grille et les composants.
  \item \textbf{Chart.js} chargé via CDN dans \texttt{App.razor}, scripts personnalisés \texttt{portfolio-charts.js} et \texttt{theme.js}.
  \item \textbf{Interop JavaScript} (\texttt{IJSRuntime}) depuis les composants interactifs pour initialiser et mettre à jour les graphiques.
\end{itemize}

\section{Mode de rendu}

Les pages à interactions fortes utilisent \texttt{@rendermode InteractiveServer}. Le reconnect modal Blazor est conservé pour une meilleure expérience en cas de perte de circuit SignalR.

\chapter{Validation et déploiement local}
\label{ch:val}

\section{Vérifications automatisées}

La commande \texttt{dotnet build} sur la solution ou le projet Web doit s'exécuter sans erreur. Les migrations sont appliquées au lancement ; en cas de doute, \texttt{dotnet ef database update} peut être exécuté avec le projet d'infrastructure et le projet Web comme startup.

\section{Vérifications manuelles recommandées}

\begin{enumerate}[leftmargin=*]
  \item Lancer l'application (\texttt{dotnet run} ou \texttt{dotnet watch} depuis \texttt{GestionPortefeuille.Web}).
  \item Parcourir le CRUD actifs : création, édition, suppression, filtres.
  \item Enregistrer des transactions achat/vente : vérifier la diminution/augmentation de la liquidité et le refus des opérations invalides.
  \item Consulter le tableau de bord : cohérence des pourcentages de répartition et des gains/pertes avec le CMP.
  \item Onglet Outils : import CSV de prix, backtest sur un actif disposant d'historique.
  \item Sous Windows : si \texttt{dotnet build} échoue sur des fichiers verrouillés (\texttt{.dll}, \texttt{.pdb}), arrêter toute instance en cours (\texttt{Ctrl+C}) ou terminer le processus hôte avant de recompiler.
\end{enumerate}

\section{Limites connues}

\begin{itemize}[leftmargin=*]
  \item Historique de prix \textbf{simulé} : les conclusions du backtest et de la volatilité sont illustratives.
  \item APIs marché soumises à quotas, clés et disponibilité réseau ; le mapping crypto est limité à une liste de symboles.
  \item Pas d'authentification multi-utilisateur dans la version décrite (évolution possible).
\end{itemize}

\chapter{Conclusion et perspectives}
\label{ch:conc}

\section{Bilan}

Le projet livre une application Web \textbf{complète} répondant aux exigences du module : Blazor, EF Core Code First avec migrations, architecture en couches, CRUD et recherche LINQ, tableau de bord avec graphiques, validation par annotations, code asynchrone. L'extension \textbf{data science} renforce la dimension analytique (risque, tendance, backtest) et supporte la partie « créativité » du barème.

\section{Pistes d'évolution}

\begin{itemize}[leftmargin=*]
  \item \textbf{Authentification} (ASP.NET Core Identity) et portefeuilles utilisateur isolés.
  \item \textbf{Tests unitaires} sur \texttt{PositionCostHelper}, règles de transaction et calculs d'analytique.
  \item \textbf{CI/CD} (GitHub Actions) : build + tests + publication.
  \item \textbf{API REST} optionnelle pour alimenter une mobile app ou un tableau externe.
  \item \textbf{Coût moyen} plus avancé (FIFO réel) et frais de transaction dans le modèle.
  \item Base \textbf{PostgreSQL} ou \textbf{SQL Server} pour un déploiement multi-utilisateurs.
\end{itemize}

\section{Remerciements / répartition travail d'équipe}

\textit{[À compléter : rôles des membres, répartition des tâches, difficultés rencontrées.]}


\appendix
\chapter{Annexe A — Extraits de code significatifs}
\label{ann:code}

\section{Point d'entrée et migrations (\texttt{Program.cs})}

\begin{lstlisting}[language=CSharp, caption={\texttt{GestionPortefeuille.Web/Program.cs} (structure principale)}]
var builder = WebApplication.CreateBuilder(args);

builder.Services.AddRazorComponents()
    .AddInteractiveServerComponents();
builder.Services.AddDbContext<AppDbContext>(options =>
    options.UseSqlite(builder.Configuration.GetConnectionString("DefaultConnection")));
builder.Services.AddPortfolioServices(builder.Configuration);

var app = builder.Build();
// ... pipeline HTTP (fichiers statiques, antiforgery, Razor Components) ...

using (var scope = app.Services.CreateScope())
{
    var db = scope.ServiceProvider.GetRequiredService<AppDbContext>();
    db.Database.Migrate();
    await DbInitializer.SeedAsync(db);
}

app.Run();
\end{lstlisting}

\section{Injection de dépendances (\texttt{AddPortfolioServices})}

\begin{lstlisting}[language=CSharp, caption={\texttt{GestionPortefeuille.Services/ServiceCollectionExtensions.cs}}]
public static IServiceCollection AddPortfolioServices(
    this IServiceCollection services, IConfiguration configuration)
{
    services.Configure<AlertOptions>(configuration.GetSection(AlertOptions.SectionName));
    services.Configure<PriceDataOptions>(configuration.GetSection(PriceDataOptions.SectionName));

    services.AddMemoryCache();
    services.AddHttpClient<IPriceDataService, PriceDataService>(client =>
    {
        client.DefaultRequestHeaders.UserAgent.ParseAdd("GestionPortefeuille/1.0 (demo)");
        client.Timeout = TimeSpan.FromSeconds(20);
    });

    services.AddScoped<IAssetService, AssetService>();
    services.AddScoped<ITransactionService, TransactionService>();
    services.AddScoped<IPortfolioService, PortfolioService>();
    services.AddScoped<IAnalyticsService, AnalyticsService>();
    services.AddScoped<IPortfolioAlertService, PortfolioAlertService>();
    return services;
}
\end{lstlisting}

\section{Lecture complémentaire}

Les entités et interfaces se trouvent dans \texttt{GestionPortefeuille.Core}, le \texttt{DbContext} et les migrations dans \texttt{GestionPortefeuille.Infrastructure}, les pages Blazor sous \texttt{GestionPortefeuille.Web/Components}. Un diagramme de classes Mermaid est fourni à la racine du dépôt dans \texttt{DiagrammeClasses.md}.

\chapter{Annexe B — Commandes utiles (.NET CLI)}
\label{ann:cmd}

\section{Restaurer et compiler}

\begin{verbatim}
cd "chemin\vers\GestionPortefeuille2"
dotnet restore
dotnet build GestionPortefeuille.sln
\end{verbatim}

\section{Exécuter l'application Web}

\begin{verbatim}
cd GestionPortefeuille.Web
dotnet run
\end{verbatim}

Mode développement avec rechargement à chaud :

\begin{verbatim}
dotnet watch run
\end{verbatim}

\section{Migrations Entity Framework Core}

Depuis la racine de la solution (adapter les chemins si besoin) :

\begin{verbatim}
dotnet ef migrations add NomMigration ^
  --project GestionPortefeuille.Infrastructure ^
  --startup-project GestionPortefeuille.Web

dotnet ef database update ^
  --project GestionPortefeuille.Infrastructure ^
  --startup-project GestionPortefeuille.Web
\end{verbatim}

\textit{Remarque :} sous PowerShell, remplacer \verb|^| par le caractère d'échappement de continuation de ligne adéquat ou écrire la commande sur une seule ligne.

\section{Problème de fichiers verrouillés (Windows)}

Si le build échoue avec un message indiquant qu'un \texttt{.dll} ou \texttt{.pdb} est utilisé par un autre processus :

\begin{enumerate}[leftmargin=*]
  \item Arrêter l'application (\texttt{Ctrl+C} dans le terminal où \texttt{dotnet watch} ou \texttt{dotnet run} tourne).
  \item Si nécessaire : \texttt{dotnet build-server shutdown}, puis relancer \texttt{dotnet build}.
\end{enumerate}

\section{Compiler le PDF du rapport}

\begin{verbatim}
cd Rapport
pdflatex -interaction=nonstopmode rapport.tex
pdflatex -interaction=nonstopmode rapport.tex
\end{verbatim}

Ou avec \texttt{latexmk} :

\begin{verbatim}
latexmk -pdf -interaction=nonstopmode rapport.tex
\end{verbatim}


\end{document}
